\def \titolo {Meccanica Quantistica}

\documentclass[11pt,twoside]{scrbook} % KOMA

% ---------- LINGUA & CODIFICA ----------
\usepackage[italian]{babel}
\usepackage[T1]{fontenc}

% ---------- TIPOGRAFIA ----------
\usepackage{lmodern}
\usepackage{microtype}

% ---------- GEOMETRIA PAGINA ----------
\usepackage[
  a4paper,
  top=24mm,
  bottom=28mm,
  inner=24mm,   % margine interno (verso rilegatura)
  outer=30mm,   % margine esterno: note a margine
  heightrounded
]{geometry}
\setlength{\headheight}{17pt} % evita sovrapposizioni con l'header
\setlength{\headsep}{6mm}     % distanza tra header e corpo

\usepackage[automark]{scrlayer-scrpage}
\clearpairofpagestyles
\automark[section]{chapter}
\ohead{\headmark}
\ofoot*{\pagemark}
\pagestyle{scrheadings}

% ---------- STRUMENTI DI STRUTTURA ----------
\setcounter{secnumdepth}{3}
\setcounter{tocdepth}{3}

% ---------- STILE TITOLI IN STILE KOMA ----------
\setkomafont{chapter}{\normalfont\huge\bfseries}
\setkomafont{section}{\normalfont\Large\bfseries}
\setkomafont{subsection}{\normalfont\large\bfseries}
\setkomafont{subsubsection}{\normalfont\normalsize\bfseries}

% ---------- STILE FRONTESPIZIO ----------
\setkomafont{title}{\sffamily\bfseries\Huge}
\setkomafont{author}{\sffamily\large}
\setkomafont{publishers}{\sffamily\normalsize}
\setkomafont{date}{\large}

\RedeclareSectionCommand[beforeskip=1.5ex,afterskip=1.5ex]{chapter}
\RedeclareSectionCommand[beforeskip=1ex,afterskip=0.8ex]{section}
\RedeclareSectionCommand[beforeskip=0.8ex,afterskip=0.6ex]{subsection}
\RedeclareSectionCommand[beforeskip=0.6ex,afterskip=0.4ex]{subsubsection}

% ---------- AMBIENTI UTILI ----------
\usepackage{amsmath,amssymb,amsthm}
\usepackage{booktabs}
\usepackage{array}
\usepackage{graphicx}
\usepackage{xcolor}
\usepackage{tcolorbox}
\usepackage{subcaption} 
\usepackage{upgreek} % \uptau
\usepackage{mathrsfs} % \mathscr{}
\usepackage{dsfont}
\usepackage{tikz}
\usetikzlibrary{intersections}
\usepackage{pgfplots} % grafici, creato in tikZ
\pgfplotsset{compat = 1.18}
\usetikzlibrary{
arrows.meta,
positioning,
calc,
decorations.markings,
angles,
quotes,
3d,
patterns
}
\usepackage[european]{circuitikz} % circuiti
\usepackage{comment} % \begin{comment}  \end{comment}
\usepackage{mathtools} % \coloneqq
\usepackage{siunitx}
\usepackage{hyperref}
\hypersetup{  % stile dei riferimenti
  colorlinks=true,
  linkcolor=black,
  citecolor=blue,
  urlcolor=blue
}
\sisetup{separate-uncertainty=true}
%----------- TEOREMI E STRUTTURE VARIE (KOMA, AMSTHM)

% Stile "teorema" (corsivo)
\theoremstyle{plain}
\newtheorem{theorem}{Teorema}[chapter]
\newtheorem{proposition}[theorem]{Proposizione}
\newtheorem{lemma}[theorem]{Lemma}
\newtheorem{corollary}[theorem]{Corollario}

% Stile "definizione" (testo dritto)
\theoremstyle{definition}
\newtheorem{definition}[theorem]{Definizione}

% --- Postulati: senza numero o con numero scelto a mano ---
\newtheorem*{postulate*}{Postulato}       % senza numerazione
\newtheorem{postulateinner}{Postulato}    % d'appoggio per il numero manuale
\NewDocumentEnvironment{postulate}{o o}
  {\IfNoValueTF{#1}
     {\begin{postulate*}}
     {\renewcommand{\thepostulateinner}{#1}%
      \IfNoValueTF{#2}{\begin{postulateinner}}{\begin{postulateinner}[#2]}}}
  {\IfNoValueTF{#1}{\end{postulate*}}{\end{postulateinner}}}

% Stile "nota" (più discorsivo)
\theoremstyle{remark}
\newtheorem{remark}[theorem]{Osservazione}
\newtheorem{example}[theorem]{Esempio}

% ---------- NOTE A MARGINE ----------
% \marginpar stampa nel margine esterno con oneside.
\newcommand{\mkeyword}[1]{%
  \marginpar{\raggedright\footnotesize\emph{#1}}%
}

% per le note a sinistra:
% \reversemarginpar

% ---------- PARAGRAFI E SPAZI ----------
\setlength{\parindent}{0.8em}
\setlength{\parskip}{0.25em}

%----------- COMANDI AGGIUNTIVI -----------

\newcommand{\1}{\mathds{1}} % in stile funzione caratteristica

% --- vettore colonna ---
\newcommand{\colvec}[1]{%
  \begin{pmatrix}
    #1
  \end{pmatrix}%
}

% --- derivata parziale ---
\newcommand{\der}[3][]{%
  \frac{\partial^{#1} {#2}}{\partial {#3}^{#1}}%
}

% --- versore --- 
\newcommand{\vers}[1]{
  \hat{\boldsymbol{#1}}
}

% --- epsilon ---
\newcommand{\eps}[0]{
  \varepsilon
}

% --- norma ---
\newcommand{\norm}[1]{
  \left\lVert #1 \right\rVert
}

% --- prodotto scalare ---
\newcommand{\sc}[1]{
  \langle #1 \rangle
}

% --- prodotto scalare aggiustabile ---
\newcommand{\sca}[1]{
  \left< #1 \right>
}

% --- vettore ---
\renewcommand{\v}[1]{
  \mathbf{#1}
}

% --- \backslash ---
\newcommand{\bs}[0]{
	\backslash 
}

% --- \diag ---
\newcommand{\diag}[0]{
	\textnormal{diag} 
}

% --- Logaritmo principale ---
\newcommand{\Log}[0]{
  \textnormal{Log }
}

% --- overline ---
\newcommand{\ol}[1]{
  \overline{#1}
}

% --- underline ---
\newcommand{\ul}[1]{
  \underline{#1}
}

% --- valore assoluto ---
\newcommand{\abs}[1]{
  \left|{#1}\right|
}

% --- non andare a capo --- NON FUNZIONA
\newcommand{\inline}[1]{
  \binoppenalty=10000 \relpenalty=10000 #1
}

% --- Residui ---
\DeclareMathOperator{\Res}{
  \textnormal{Res}
}

% --- Supporto ---
\DeclareMathOperator{\supp}{
  \textnormal{supp}
}

% --- Kernel ---
\DeclareMathOperator{\Ker}{
  \textnormal{Ker}
}

% --- Rango ---
\DeclareMathOperator{\rank}{
  \textnormal{rank}
}

% % --- Kernel ---
% \DeclareMathOperator{\Ker}{
%   \textnormal{Ker}
% }

% --- Overset ---
\newcommand{\ov}[2]{
  \overset{#2}{#1}
}

% --- Underset ---
\newcommand{\un}[2]{
  \underset{#2}{#1}
}

% --- Parte principale ---
\let\P\relax
\DeclareMathOperator{\P}{
  \mathcal{P}
}

% --- Funzione segno ---
\DeclareMathOperator{\sgn}{
  \textnormal{sgn}
}

% --- Traccia ---
\DeclareMathOperator{\Tr}{
  \textnormal{Tr}
}

% --- Probabilità ---
\DeclareMathOperator{\Prob}{
  \textnormal{Prob}
}

% =================== bra-ket ===================

% --- bra ---
\renewcommand{\b}[1]{
  \langle #1 \vert
}

% --- bra-2 ---
\newcommand{\ba}[1]{
  \left< #1 \right|
}

% --- ket ---
\renewcommand{\k}[1]{
  \vert #1 \rangle
}

% --- ket-2 ---
\newcommand{\ka}[1]{
  \left| #1 \right>
}

% % --- bracket ---
% \newcommand{\bk}[2]{
%   \left< #1 \vert #2 \right>
% }

% --- bracket ---
\newcommand{\bk}[1]{
  \langle #1 \rangle
}

% --- bracket-aggiustabile ---
\newcommand{\bka}[1]{
  \left< #1 \right>
}

% =================== Spazi ===================

% --- R numeri reali --- 
\newcommand{\R}[0]{
  \mathbb{R}
}

% --- C numeri complessi --- 
\newcommand{\C}[0]{
  \mathbb{C}
}

% --- E Funzioni lisce --- 
\newcommand{\E}[0]{
  \mathcal{E}
}

% --- S funzioni a decrescenza rapida --- 
\let\S\relax
\newcommand{\S}[0]{
  \mathcal{S}
}

% --- D Funzioni lisce a supporto compatto --- 
\newcommand{\D}[0]{
  \mathcal{D}
}

% =================== Trasformate ===================

% --- Fourier ---
\DeclareMathOperator{\F}{
  \mathcal{F}
}

% --- \hat{} ---
\newcommand{\h}[1]{
  \hat{#1}
}

% --- Laplace ---
\let\L\relax
\DeclareMathOperator{\L}{
  \mathcal{L}
}

% --- \tilde{} ---
\newcommand{\ti}[1]{
  \tilde{#1}
}

% --- Hilbert ---
\let\H\relax
\DeclareMathOperator{\H}{
  \mathcal{H}
}



% =================== Scambi ===================

% --- phi φ ϕ ---
\let \ORIGINALphi \phi
\let \ORIGINALvarphi \varphi

\renewcommand{\phi}{\ORIGINALvarphi}
\renewcommand{\varphi}{\ORIGINALphi}

% =================== Stili ===================

% --- displaystyle ---
\newcommand{\DS}[0]{
	\displaystyle 
}

% --- textstyle ---
\newcommand{\TS}[0]{
	\textstyle 
}

% % --- scriptstyle --- nota: non attivi per mantere liberi i comandi \SS \SSS
% \renewcommand{\SS}[0]{
% 	\scriptstyle 
% }

% % --- \displaystyle ---
% \newcommand{\SSS}[0]{
% 	\scriptscriptstyle 
% }

% =================== Template ===================

% \mathrm{F}, \mathit{F}, \mathbf{F}, \mathsf{F}, \mathtt{F}, \boldsymbol{F}, \mathcal{F}, \mathscr{F}, \mathbb{F}, \mathds{F}, \mathfrak{F}

% ---------- BIBLIOGRAFIA ----------
\usepackage{csquotes} % \enquote{}
\usepackage[backend=biber,style=numeric,sorting=nyt]{biblatex}
\addbibresource{bibliografia.bib} % nome del file .bib

% ---------- FRONTESPIZIO ----------
\title{\textbf{\titolo}}
\author{Cosimo Pascarella}
\date{Anno Accademico 2026--2027}

\publishers{
  Email: \text{c.pascarella2@studenti.unipi.it}
}

\begin{document}
\frontmatter
\maketitle
\tableofcontents
\cleardoublepage
\markboth{}{}
\mainmatter

\chapter{Prerequisiti matematici}
Dato un sistema ortonormale $ \{\k{\psi_n}\} $, la traccia di un operatore $ A $ è definita come 
\begin{align}
  \Tr A = \sum_{n} \b{\psi_n} A \k{\psi_n}
\end{align}
ed è indipendente dal sistema ortonormale scelto. Tale definizione nasce dal caso finito e si può estendere senza troppi problemi anche al caso infinito, 
tuttavia in alcuni casi dovremo preoccuparci se la somma diverge o meno.


\chapter{Formulazione della meccanica quantistica}

\begin{postulate}[1]
  A ogni variabile dinamica corrisponde un operatore lineare, e i possibili valori di tale variabile sono gli autovalori dell'operatore lineare associato.
\end{postulate}

\begin{postulate}[ 2 ]
  A ogni stato corrisponde un unico \textit{operatore di stato}. Il valor medio di una variabile dinamica $\mathbf{R}$, rappresentata dall'operatore $ R $, 
  nell'ensemble virtuale di eventi risultanti da una procedura di preparazione per lo stato, rappresentato dall'operatore $ \rho $, vale 
  \begin{align}
    \sca{\mathbf{R}} = \frac{\Tr(\rho R)}{\Tr \rho}.
  \end{align}
\end{postulate}

Secondo l'interpretazione statistica, la teoria non predice il singolo esito ma le distribuzioni statistiche per una serie di preparazioni identiche. 
Con questa interpretazione possiamo definire uno stato come ciò che fissa un determinato esito, quindi diremo che due preparazioni sono nello stesso stato 
se per ogni variabile dinamica danno le stesse statistiche. In questo senso il postulato 2 ci dice che l'informazione che codifica uno stato è rappresentata 
da un operatore lineare $ \rho $, ed è definito tramite la proprietà che applicato a un operatore $ R $ associato a una variabile dinamica, presa la traccia 
ne restituisce la media sull'ensemble (moltiplicata per la traccia di $ \rho $).

Possiamo imporre alcuni vincoli sulla media $ \sca{\mathbf{R}} $ in modo da ricavare delle proprietà dell'operatore di stato $ \rho $. 
\begin{enumerate}
  \item \textbf{Normalizzazione}: gli operatori $ \rho $ e $ c \rho $ con $ c \in \C $ danno le stesse medie dato che $ \Tr (c\rho) = c \Tr \rho $, per 
  questo motivo possiamo imporre che $ \Tr \rho = 1 $, che corrisponde a scegliere $ \DS c = \frac1{\Tr\rho} $.
  \item \textbf{Hermiticità}: consideriamo un osservabile a cui è associato un operatore $ P_u = \k{u} \b{u} $ che rappresenta un proiettore 
  su un vettore $ \k{u} $ di norma 1. Dato che la media di una variabile dinamica deve essere reale (perché assume solo valori 
  reali) allora lo deve essere anche $ \Tr (\rho P_u) = \sum_n \b{\psi_n} \rho \k{u} \bk{u|\psi_n} = \b{u} \rho \k{u} $, e siccome deve valere per ogni 
  $ \k{u} $ allora si ha che $ \rho = \rho^\dagger $ per il lemma \ref{lem:}.
  \item \textbf{Positività}: considerando ancora i proiettori, gli autovalori possibili sono $ 0 $ e $ 1 $, per tanto la media dell'osservabile 
  deve essere positiva e quindi $ 0 \leq \sca{\mathbf{R}} = \b{u} \rho \k{u} $. Siccome deve valere per ogni $ \k{u} $ allora $ \rho \geq 0 $.
\end{enumerate}
Queste tre proprietà di un generico operatore di stato $ \rho $, come abbiamo visto, discendono solo da considerazioni sulle caratteristiche di una misura, 
possiamo quindi riformulare il secondo postulato con il seguente.
\begin{postulate}[ 2a ]
  A ogni stato corrisponde un unico operatore di stato di traccia unitaria, hermitiano e non negativo.
\end{postulate}
Dove l'unicità è dovuta al fatto che dicevamo sopra cioè che se due preparazioni danno le stesse statistiche allora sono lo stesso stato. A questo punto 
possiamo scambiare i ruoli e consideriamo come operatore di stato $ \rho = \k{\psi} \b{\psi} $ con $ \k{\psi} $ di norma unitaria, 
il quale rispetta le tre condizioni del postulato 2a. Per un operatore $ R $ associato a un osservabile $ \mathbf{R} $ abbiamo che 
\begin{align}
  \sca{\mathbf{R}} = \Tr(\rho R) = \sum_{n} \bk{\phi_n|\psi} \b{\psi} R \k{\phi_n} = \b{\psi} R \k{\psi}
\end{align}
che deve essere reale per ogni $ \k{\psi} $ e quindi si ottiene che $ R = R^\dagger $. In questo modo possiamo ridefinire anche il postulato 1 con il 
seguente.
\begin{postulate}[ 1a ]
  A ogni variabile dinamica corrisponde un operatore hermitiano i cui autovalori sono i possibili valori della variabile dinamica.
\end{postulate}

Dato che l'operatore di stato $ \rho $ è autoaggiunto per il teorema spettrale possiamo rappresentarlo come 
\begin{align}
  \rho = \sum_n \rho_n \k{\psi_n} \b{\psi_n}
\end{align}
con $ \k{\psi_n} $ la base degli autovettori e $ \rho_n $ i suoi autovalori (assunti discreti) che sono reali. Da questa rappresentazione deduciamo che 
$ 1 = \Tr \rho = \sum_n \rho_n $, sono anche tutti non negativi dato che $ \b{\psi} \rho \k{\psi} \geq 0 $ per ogni $ \k{\psi} $. Combinando queste 
proprietà si ottiene che 
\begin{align}
  0 \leq \rho_n \leq 1
\end{align}
e l'insieme di tutti gli operatori con queste caratteristiche forma un \textit{insieme convesso}, cioè che preso un insieme di operatori 
$ \{\rho^{(i)}\} $ con tali proprietà allora anche $ \rho = \sum_{i} a_i \rho^{(i)} $ rispetta tali proprietà, a patto che $ \sum_i a_i = 1 $ e 
$ 0 \leq a_i \leq 1 $.

\section{Stati puri}

\begin{definition}[  ]
  Uno \textit{stato puro} è uno stato con un operatore di stato $ \rho $ della forma 
  \begin{align}
    \rho = \k{\psi} \b{\psi} 
  \end{align}
  dove $ \k{\psi} $ ha norma unitaria ed è detto \textit{vettore di stato}.
\end{definition}

In uno stato puro il valor medio di un osservabile $ \mathbf{R} $ è 
\begin{align}
  \mathbf{R} = \Tr\left( \k{\psi} \b{\psi} R \right) = \b{\psi} R \k{\psi}
\end{align}
e il vettore di stato non è unico infatti per ogni $ \alpha \in \R $ si ha che tutti i vettori di stato $ e^{i\alpha} \k{\psi} $ danno lo stesso stato 
puro. Uno stato puro è quindi il proiettore sul vettore di stato (a meno di una fase) e quindi una caratterizzazione equivalente è che 
\begin{align}
  \rho^2 = \rho
\end{align}
che è necessaria per la definizione di stato puro ed è anche sufficiente perché gli autovalori $ \rho_n $ devono soddisfare $ \rho_n^2 = \rho_n $ quindi 
possono essere solo $ 0,1 $. Per le caratteristiche degli operatori di stato deve essere che la somma degli autovalori è $ 1 $ e quindi l'unica possibilità 
che $ \rho_\psi = 1 $ e tutti gli altri $ 0 $, con rappresentazione spettrale $ \rho = \sum_n \rho_n \k{\psi_n} \b{\psi_n} = \k{\psi} \b{\psi} $. 
Una terza condizione equivalente è che 
\begin{align}
  \Tr \left( \rho^2 \right) = 1
\end{align}
che è necessaria per quello che abbiamo appena detto, per la sufficienza dato che $ 0 \leq \rho_n \leq 1 $ si ha $ \rho_n^2 \leq \rho_n $ ma allora 
$ \Tr(\rho^2) \leq \Tr(\rho) = 1 $, abbiamo l'uguaglianza se $ \rho_n^2 = \rho_n $ e quindi è equivalente alla seconda caratterizzazione. Un'ulteriore 
condizione è data dal seguente teorema.

\begin{theorem}[  ]
  Uno stato puro non può essere mai espresso come una combinazione convessa non banale di altri stati, mentre uno stato non puro può esserlo.
\end{theorem}

Il precedente teorema ci dice che gli unici stati che si possono esprimere come combinazione convessa sono gli stati non puri, mentre per gli stati puri 
vale solo quella banale e cioè con se stessi. Ad esempio $ \rho = \lambda \rho_1 + (1 - \lambda) \rho_2 $ con $ \lambda \in (0,1) $ è necessariamente 
uno stato non puro.

\section{Distribuzioni di probabilità}

Vogliamo ottenere la distribuzione dei valori dell'osservabile $ \mathbf{R} $ sull'ensemble, per farlo sfruttiamo una funzione ausiliaria $ F(\mathbf{R}) $ 
con l'operatore associato $ F(R) $ costruito come
\begin{align}
  F(R) = \sum_{n} F(r_n) \k{\psi_n} \b{\psi_n}
\end{align}
in modo che gli autovettori siano conservati. Sia $ g(r)dr $ la probabilità che $ \mathbf{R} $ sia fra $ r $ e $ r + dr $, allora scriveremo che 
\begin{align}
  \sca{\mathbf{R}} &= \int_{-\infty}^{+\infty} r \, g(r) \,dr
  \\
  \sca{F(\mathbf{R})} &= \int_{-\infty}^{+\infty} F(r) \,g(r) \,dr
\end{align}
ma per il postulato 2 vale che $ \sca{F(\mathbf{R})} = \Tr(\rho F(R)) $ quindi scegliendo una opportuna $ F(R) $ possiamo ottenere la densità di 
probabilità $ g(r) $. 

\textbf{Spettro discreto}: sia $ R $ un operatore autoaggiunto con spettro discreto, possiamo rappresentarlo con i suoi autovalori $ r_n $ e in un sistema 
di autovettori ortonormali $ \k{r_n} $ come 
\begin{align}
  R = \sum_n r_n \k{r_n} \b{r_n}.
\end{align}
Come funzione ausiliaria usiamo $ F(\mathbf{R}) = \theta(r - \mathbf{R}) $ e otteniamo 
\begin{align}
  \sca{\theta(r - \mathbf{R})} = \int_{-\infty}^{r} g(r') \,dr' = \Prob (\mathbf{R} < r \vert \rho)
\end{align}
mentre per il postulato 2 si ha che
\begin{align}
  \sca{\theta(r - \mathbf{R})} = \Tr\left[ \rho \,\theta(r - R) \right] = \Tr \left[ \rho \sum_n \theta(r - r_n) \k{r_n} \b{r_n} \right] 
  = \sum_n \b{r_n} \rho \k{r_n} \theta(r - r_n)
\end{align}
dove abbiamo usato che il sistema $ \k{r_n} $ è ortonormale. Combinando le due espressioni si ottiene che 
\begin{align}
  g(r) &= \der{}{r} \Prob (\mathbf{R} < r \vert \rho) = \der{}{r} \sca{\theta(r - \mathbf{R})} = \der{}{r} \sum_n \b{r_n} \rho \k{r_n} \theta(r - r_n)
  \\
  &= \sum_n \b{r_n} \rho \k{r_n} \der{}{r} \theta(r - r_n) = \sum_n \b{r_n} \rho \k{r_n} \delta(r - r_n) 
\end{align}
dove si nota chiaramente che la densità di probabilità è $ 0 $ se $ r $ non è un autovalore, consistente con il postulato 1. Calcoliamo quindi la 
probabilità che $ \mathbf{R} $ assuma il valore $ r $ nello stato $ \rho $:
\begin{align}
  \Prob(\mathbf{R} = r \vert \rho) &= \lim_{\eps \to 0} \left\{\Prob(\mathbf{R} < r + \eps \vert \rho) - \Prob(\mathbf{R} < r - \eps \vert \rho) \right\}
  \\
  &= \int_{r-\eps}^{r + \eps} g(r')\,dr' = \sum_n \b{r_n} \rho \k{r_n} \delta_{r,r_n}
\end{align}
dove $ \delta_{r,r_n} $ è la delta di Kronecker. Nel caso di uno stato puro $ \rho = \k{\psi} \b{\psi} $ e autovalori non degeneri la probabilità diventa 
\begin{align}
  \Prob(\mathbf{R} = r_n \vert \rho) = \bk{r_n|\psi} \bk{\psi|r_n} = \abs{\bk{r_n|\psi}}^2.
\end{align}
Se $ \mathbf{R} $ assume l'autovalore non degenere $ r_0 $ con probabilità $ 1 $ allora deve valere che $ \b{r_0} \rho \k{r_0} = 1 $, 
ma per un qualsiasi stato vale che 
\begin{align}
  1 \geq \Tr\left( \rho^2 \right) = \sum_{m} \b{r_m} \rho^2 \k{r_m} = \sum_{n,m} \b{r_m} \rho \k{r_n} \b{r_n} \rho \k{r_m} 
  = \sum_{n,m} \abs{\b{r_m} \rho \k{r_n}}^2
\end{align}
e tale somma è saturata dal termine $ \b{r_0} \rho \k{r_0} = 1 $ quindi $ \b{r_m} \rho \k{r_n} = 0 $ per ogni $ (m,n) \neq (0,0) $, ma allora l'unico stato 
in cui $ \mathbf{R} $ assume $ r_0 $ è lo stato puro $ \rho = \k{r_0} \b{r_0} $. Uno stato di questo tipo è detto \textit{autostato} 
dell'osservabile $ \mathbf{R} $.

\textbf{Spettro continuo}: sia $ Q $ un operatore autoaggiunto associato all'osservabile $ \mathbf{Q} $ con spettro solo continuo, possiamo quindi 
rappresentarlo come 
\begin{align}
  Q = \int q \k{q} \b{q} \,dq
\end{align}
dove i suoi autovettori $ \k{q} $ soddisfano la relazione di completezza $ \int \k{q} \b{q} \,dq = 1 $ e la relazione di ortogonalità 
$ \bk{q'|q''} = \delta(q' - q'') $. Come per lo spettro discreto abbiamo che per $ F(\mathbf{Q}) = \theta(q - \mathbf{Q}) $ si ottiene 
\begin{align}
  \sca{\theta(r - \mathbf{Q})} = \int_{-\infty}^{r} g(q') \,dq' = \Prob (\mathbf{Q} < q \vert \rho)
\end{align}
che è la probabilità che l'osservabile $ \mathbf{Q} $ sia minore di $ q $ nello stato $ \rho $. Tuttavia per il secondo postulato ora abbiamo 
\begin{align}
  \sca{\theta(r - \mathbf{Q})} = \Tr \left[ \rho \, \theta(q - Q) \right] = \Tr\left[ \rho \int_{-\infty}^{q} q'\k{q'} \b{q'} \,dq' \right]
  = \int_{-\infty}^{q} \b{q'}\rho \k{q'}\,dq'
\end{align}
e quindi la densità di probabilità $ g(q) $ vale 
\begin{align}
  g(q) = \der{}{q} \Prob(\mathbf{Q} < q | \rho) = \b{q} \rho \k{q}
\end{align}
che nel caso di uno stato puro $ \rho = \k{\psi} \b{\psi} $ diventa
\begin{align}
  g(q) = \bk{q | \psi} \bk{\psi | q} = \abs{\bk{q | \psi}}^2.
\end{align}

\chapter{Cinematica e dinamica}



\section{Trasformazioni}
Vogliamo capire come cambiano gli osservabili e gli stati (puri) a seguito di trasformazioni nello spazio e nel tempo. La fisica del sistema non 
deve cambiare sotto rotazioni, traslazioni o se viene messo in moto a velocità costante, perciò deve esistere una relazione che le lega lo stato 
$ k{\psi} $ allo stato trasformato $ \k{\psi'} $ e l'osservabile $ A $ con il trasformato $ A' $. Inoltre, devono essere conservate le distribuzioni 
di probabilità $ \abs{ \bk{\phi | \psi}}^2 $, e le trasformazioni che hanno questa proprietà sono descritte dal seguente teorema.

\begin{theorem}[ di Wigner ]
  Ogni trasformazione di uno spazio vettoriale $ \H $ in sé stesso che preserva il modulo del prodotto scalare $ \abs{\bk{\phi | \psi}} $ può essere 
  rappresentata con un operatore $ U $ o unitario o antiunitario.
\end{theorem}

Nel caso che $ U $ sia unitario mostriamo che preserva effettivamente il prodotto scalare (non solo il suo modulo):
\begin{align}
  \bk{\phi' | \psi'} = \bk{U \phi | U \psi} = (\b{\phi} U^\dagger) \, (\k{U |\psi}) = \bk{\phi | U^\dagger U | \psi} = \bk{\phi | \psi}
\end{align}
dato che $ U^\dagger U = I $. Mentre se $ U $ è antiunitario, cioè antilineare e unitario, la mappa $\psi \mapsto \bk{\phi | U \psi}$ 
è antilineare in $ \psi $ perché 
\begin{align}
  \bk{\phi | U(\alpha\, \psi)} = \alpha^* \bk{\phi | U\psi}
\end{align}
e quindi non è un funzionale lineare su $ \H $, perciò dobbiamo definire l'applicazione di $ U $ a $ \b{\phi} $ come 
\begin{align}
  (\b{\phi} U) \k{\psi} \coloneq \bk{\phi | U\psi }^*
\end{align}
che è un elemento del duale $ \H^* $. Dunque abbiamo che 
\begin{align}
  \bk{U \phi | U \psi} = (\b{\phi} U^\dagger) (U \k{\psi}) = {\bk{\phi | U^\dagger U | \psi}}^* = \bk{\phi |\psi}^*.
\end{align}
Gli operatori antilineari possono descrivere trasformazioni discrete ma non trasformazioni continue, questo perché una trasformazione $ U(l) $ 
con $ l $ parametro continuo può essere vista come una combinazione di altre trasformazioni ad esempio $ U(l) = U(l/2) U(l/2) $ e il prodotto 
di due operatori antilineari è lineare perché si coniuga due volte, quindi a prescindere dal carattere di $ U(l/2) $ si ha che $ U(l) $ è lineare.

Nel trasformare gli operatori si devono conservare gli autovalori e gli autovettori così da avere gli stessi valori possibili delle misure, infatti 
abbiamo
\begin{align}
  A'\k{\phi'_n} = a_n \k{\phi'_n} 
\end{align}
e usando che $ \k{\phi'_n} = U \k{\phi_n} $ si ha
\begin{align}
  A' U \k{\phi_n } = a_n U\k{\phi_n} \implies U^{-1} A' U \k{\phi_n} = a_n \k{\phi_n}.
\end{align}
A questo punto per concludere che $ A = U^{-1} A' U $ ci basta osservare che $ (A - U^{-1} A' U )\k{\phi_n} = 0 $ vale per ogni vettore 
del set completo $ \{\k{\phi_n}\} $ e quindi vale per ogni vettore. Invertendo si ottiene  
\begin{align}
  A = U A' U^{-1}.
\end{align}

Consideriamo una famiglia di operatori unitari $ U(s) $ che dipendono dal parametro continuo $ s $,





\end{document}